A força de tração no cabo é dada pela equação:
\begin{equation}
  T = \frac{W_s}{n_c \cdot \eta_{\text{moitão}}} \cdot 10^{-1}
  \label{eq:forca_de_tracao_no_cabo}
\end{equation}
Onde:
\begin{itemize}[leftmargin=1.25cm]
  \item $T$: força de tração no cabo em $\si{\deca\newton}$.
  \item $W_s$: carga útil somada ao peso do moitão em $\si{\newton}$.
  \item $n_c$: número de cabos.
  \item $\eta_{\text{moitão}}$: rendimento do moitão.
\end{itemize}
\vspace{1em}

A carga útil somada ao peso do moitão é dada pela equação:
\begin{equation}
  W_s = W_u + G_{\text{moitão}}
  \label{eq:ws}
\end{equation}
Onde:
\begin{itemize}[leftmargin=1.25cm]
  \item $W_s$: carga útil somada ao peso do moitão em $\si{\newton}$.
  \item $W_u$: carga útil em $\si{\newton}$.
  \item $G_{\text{moitão}}$: peso do moitão em $\si{\newton}$.
\end{itemize}
\vspace{1em}

Substituindo os valores na equação \eqref{eq:ws}:
\begin{align*}
  W_s & = 100 \cdot 10^3 + 750   \\
  W_s & = \SI{100,75e3}{\newton}
\end{align*}
\vspace{1em}

Substituindo os valores na equação \eqref{eq:forca_de_tracao_no_cabo}:
\begin{align*}
  T & = \frac{100,75 \cdot 10^3}{4 \cdot 0,9925} \cdot 10^{-1} \\
  T & = \boxed{\SI{2537,78}{\deca\newton}}
\end{align*}
